\documentclass[11pt]{article}

\usepackage[utf8]{inputenc}
\usepackage[T1]{fontenc}
\usepackage{amsmath,amssymb,amsthm}
\usepackage{booktabs}
\usepackage{hyperref}
\usepackage[margin=1in]{geometry}
\usepackage{url}
\def\UrlBreaks{\do\/\do-\do_}

\hypersetup{colorlinks=true, linkcolor=blue, citecolor=blue, urlcolor=blue}
\setlength{\tabcolsep}{4pt}
\newtheorem{lemma}{Lemma}
\newtheorem{fact}{Fact}

\title{Folding a Ring onto Heavy-Hex:\\
Formula-Built Schedules with Provable Gate Counts, and What Their Signal Shows on Twins and Devices}

\author{TN Holdings\\
\texttt{love.os.architect@proton.me}\\
\texttt{https://github.com/TN-Holdings-LLC/psf-zero}}

\date{Version 1 --- 26 October 2026}

\begin{document}
\maketitle

\begin{abstract}
A ring of $n$ CZ gates per layer, the entangling pattern of many variational classifiers, does not fit IBM's
heavy-hex devices unless $n$ is a cycle length of the lattice, and heavy-hex has cycles of $12, 20, 24, 28, \dots$
qubits only. We compare three schedules that a compiler can build by formula, with no layout search and without
reading the day's calibration: the \emph{zipper} (the ring folded onto a line; $5n/2-4$ two-qubit gates per layer for
even $n$, $(5n-7)/2$ for odd $n$), \emph{holes} (the ring on a longer cycle, its spare qubits moved as $|0\rangle$;
$n + 2(m-n)$), and \emph{pendants} (the ring on a shorter cycle, its extra nodes hung off it; $n + (n-m)$, on exactly
$n$ qubits). We prove that the pendant schedule's count is the least possible for layered schedules in a stated routing
model (exhaustive search finds the same for interleaved ones in small cases), that the hole schedule's count is least
within its class but not in general, and recall that no cheaper schedule than the zipper
was found by exhaustive searches at $n = 4$--$12$ and $14$. A rule that takes the fewest extra gates picks pendants at $n = 13$--$17$
and holes at $n = 8$--$11$ and $18$--$19$ on IBM's Heron devices; the rule is in PSF-Zero's AI front end a26. On noise
models, in declared matches against the zipper at 16 layers over four random classifiers, the rule won at every
$n = 13$--$18$ on two twins. Its lead was $1.5$--$8.2\times$ the zipper's kept signal on the twin of a 13-day-old
snapshot and only $1.07$--$1.28\times$ (about even at $n = 18$) on the twin of a real day's calibration, the kind that
had predicted the zipper on hardware: the old snapshot was pessimistic for the zipper. On two devices
(\texttt{ibm\_kingston}, \texttt{ibm\_fez}) the pendant schedule kept the signal its day's twin predicted
($0.90 \pm 0.12$ against $0.77$; $0.50 \pm 0.10$ against $0.59$), in two jobs of 7\,s each; these are coarse tests and
no device comparison with the zipper was made. Error-product estimates of the kept signal did not rank the schedules.
We give the benchmark number and the device user's number side by side throughout.
\end{abstract}

\section{Introduction}

Variational quantum classifiers often entangle their qubits with a ring of CZ gates in every layer
\cite{perez2020}. On a device whose coupling graph contains the ring, the layer costs $n$ CZ. IBM's Heron processors
are heavy-hex lattices \cite{chamberland2020}, which have no cycle shorter than 12 qubits and no cycle between 12 and 20
(Section~\ref{sec:cycles}). Every other ring size needs routing. A search-based compiler places and routes the ring
with the day's calibration; PSF-Zero \cite{psfzero} builds what it can by formula instead and refuses the rest. This
paper is about the formula side: which ring schedules can be written down from the coupling graph alone, how many
two-qubit gates they need, whether that is the least, and how much signal they keep.

Our earlier paper \cite{paper3} compared the zipper with Qiskit's search \cite{qiskit} through a trained classifier and
found them indistinguishable on the declared tests. Here we add two schedules that use the device's own cycles, a count
rule that chooses among the three, proofs or exhaustive checks of their counts, matches on noise models declared before
the runs, and the first device runs of the new schedule. All runs, declarations and records are in TOQB
\cite{toqb}; the compiler is PSF-Zero's AI front end a26.

\section{Heavy-hex cycles}
\label{sec:cycles}

Heavy-hex is the honeycomb lattice with every edge subdivided by one qubit \cite{chamberland2020}. A cycle in it is a
cycle of the honeycomb with every edge doubled, so its length is twice a honeycomb cycle length. Honeycomb cycles have
lengths $6, 10, 12, 14, \dots$ (one hexagon; two adjacent; three around a vertex; three in a row; \dots) and, being
bipartite with girth 6, none of length 8.
Heavy-hex cycles therefore have lengths $12, 20, 24, 28, \dots$, with none between 12 and 20. We checked this on the
coupling graphs of four Heron fake devices by listing every simple cycle up to length 22: only 12 and 20 occur (21 and
44 cycles on \texttt{FakeMarrakesh}, \texttt{FakeFez} and \texttt{FakeKingston}; 18 and 37 on \texttt{FakeTorino}).

A ring of $n$ embeds without extra gates only when $n$ is one of these lengths and a working cycle of that length exists.
Every other $n$ needs extra two-qubit gates.

\section{Three schedules built by formula}

\paragraph{Placement.} The lattice is read off the coupling graph: \emph{bridges} (qubits with two neighbours whose
indices are not consecutive), \emph{rows} (the components left without them) and \emph{cells} (between consecutive
bridges of a row pair). A cycle is the outline of a fixed shape of cells: $k$ cells side by side in a row pair
($8k+4$ qubits: 12, 20, 28), or $k$ cells plus $k_2$ cells in the next row pair down (24, 32, 36, \dots). The first
shape in a fixed order (row pairs top to bottom, cells left to right) whose qubits and couplers all work is used. No
search over mappings and no error value decides; elements at error $\ge 0.5$ are never used.

\paragraph{Zipper.} The ring laid on a line in folded order ($0, 1, 2, n-1, n-2, 3, 4, \dots$); each layer a fixed
sequence of CZ rounds, CZ+SWAP blocks at the line's ends and interior SWAPs; the next layer runs backwards. Two-qubit
gates per layer: $5n/2-4$ (even $n$), $(5n-7)/2$ (odd $n$) \cite{paper3,toqb}.

\paragraph{Holes.} The ring on the shortest working cycle of $m \ge n$ qubits; the $h = m-n$ spare qubits start in
$|0\rangle$ and are spread evenly. Each layer: CZ on every ring edge already adjacent; for each run of holes the ring
node before it moves across (CX $a \to e$, CX $e \to a$ for each empty qubit $e$), then CZ across. Two-qubit gates per layer:
$n + 2h$ when holes are apart ($n = 8$, where two holes meet, needs two more per layer in some layers).

\paragraph{Pendants.} The ring on the longest working cycle of $m \le n$ qubits; the $k = n-m$ extra ring nodes sit on
\emph{leaves}, outside neighbours of $k$ cycle nodes (\emph{hubs}), spread in cycle order. Each layer every leaf node
trades places with its hub's node by one CZ+SWAP block, which also applies their ring edge. Two-qubit gates per layer:
$n + k$, on exactly $n$ qubits; no idle qubit is carried.

\paragraph{The count rule.} Of the extra two-qubit gates per layer --- pendants $n - m_{\le}$, holes $2(m_{\ge}-n)$, the
zipper's --- the fewest decides (ties: pendants, then holes). Table~\ref{tab:rule} gives its choices on a Heron device.

\begin{table}[h]
\centering
\footnotesize
\begin{tabular}{@{}lcccccccc@{}}
\toprule
$n$ & 8--11 & 12 & 13--17 & 18--19 & 20 & 21--22 & 23 & 25--26 \\
\midrule
choice & holes (12) & ring & pendants (12) & holes (20) & ring & pendants (20) & holes (24) & pendants (24) \\
gates per layer & $16, 15, 14, 13$ & 12 & $14 \dots 22$ & $22, 21$ & 20 & $22, 24$ & 25 & $26, 28$ \\
zipper & $16, 19, 21, 24$ & 26 & $29 \dots 39$ & $41, 44$ & 46 & $49, 51$ & 54 & $59, 61$ \\
\bottomrule
\end{tabular}
\caption{The count rule's choice (cycle length in brackets) on \texttt{snapshot:marrakesh}; at $n=8$ the counts tie
and the rule takes holes. Not shown: the ring itself at $n = 24, 28$, holes at $n = 27$ and pendants at $n = 29, 30$
(cycle 28). At $n = 20$--$30$ the chosen schedule uses $0.42$--$0.47$ of the zipper's two-qubit gates.}
\label{tab:rule}
\end{table}

\paragraph{Correctness.} Every schedule equals the logical model on every measured bit: by statevector for $n \le 19$
($L = 2, 4$), by Clifford simulation with random multiples of $\pi/2$ for $n = 20$--$30$, and, for the pendant
schedule, by a structural replay of every pendant placement on cycles up to $m = 12$ (8,184 cases) and a statevector
check of every placement up to $m = 8$ (1,398 cases), done in a separate review that did not change the code. The compiler's
circuits equal TOQB's analysis arm instruction for instruction ($n = 8$--$19$).

\section{Gate-count lower bounds}

\paragraph{Routing model.} Each ring node's state sits on one qubit at a time (or on none, for an empty qubit in
$|0\rangle$). A two-qubit operation either applies a ring-edge CZ between the nodes it holds, or changes where states
sit, or both; no measurement, feed-forward, teleportation or extra entangled resource. In layer $t$ every ring edge is
applied once, after the single-qubit gates of layer $t$ on both its nodes and before those of layer $t+1$. A schedule is
\emph{layered} if all operations of layer $t$ precede those of layer $t+1$. Costs are counted in two-qubit gates; an
operation's \emph{extra} is its cost minus the ring edges it applies (at most one).

\begin{fact}
A CZ followed by a SWAP on the same pair needs two CZ (it is locally equivalent to iSWAP \cite{vatan2004,shende2004});
a SWAP needs three. A single CZ with single-qubit gates neither exchanges two states nor moves a state onto a qubit in
$|0\rangle$.
\end{fact}
\noindent For the last claim: with local gates folded in, the CZ acts on $|\chi\rangle\otimes|\varphi\rangle$ and gives
$\alpha|0\rangle|\varphi\rangle + \beta|1\rangle Z|\varphi\rangle$; if $Z|\varphi\rangle \propto |\varphi\rangle$ the
state stays on the first qubit, otherwise the output is entangled for generic $\alpha,\beta$. Hence an operation that
moves states has extra at least 1 (CZ+SWAP, which also applies a ring edge), 2 (a move onto an empty qubit, two CX) or
3 (a SWAP).

\begin{lemma}[pendants]
Let the hardware be a cycle $C_m$ plus $k < m$ leaves, each attached to a distinct hub, no two hubs adjacent, every
qubit holding a ring node ($n = m+k$). In the routing model every layer of a layered schedule costs at least $n + k$
two-qubit gates. The pendant schedule is layered and meets the bound.
\end{lemma}
\begin{proof}
Fix a layer and a leaf $\ell$ with hub $h$; let $x$ be the node on $\ell$ when the layer starts. If $x$ leaves $\ell$
during the layer, an operation that moves states acts on $(h,\ell)$, the leaf's only coupler. Otherwise both ring edges
of $x$, with $x-1$ and $x+1$, are applied across $(h,\ell)$ in this layer, so $h$ holds $x-1$ at one time and $x+1$ at
another; since no qubit is empty and $x$ stays on $\ell$, the node on $h$ changes through an operation that moves states
on a cycle edge at $h$. Name that operation $\ell$'s. Operations named for different leaves are different: one on a leaf
edge belongs to one leaf, and a cycle edge at $h$ is a cycle edge at another hub $h'$ only if $h$ and $h'$ are adjacent,
which is excluded. The layer thus holds at least $k$ distinct state-moving operations, each with extra $\ge 1$, besides
the $n$ ring edges; its cost is at least $n+k$. The pendant schedule finishes each layer before the next, with one
CZ+SWAP (extra 1) per leaf and plain CZs otherwise.
\end{proof}

On heavy-hex every coupler joins a degree-3 qubit to a degree-2 qubit and a hub has degree 3, so no two hubs are
adjacent; the coupling graphs of four Heron fake devices have no coupler between two degree-3 qubits.

\paragraph{Checks by search.} A shortest-path search over layouts and applied ring edges, from every start layout, with
the operations CZ (1), CZ+SWAP (2), SWAP (3) and move onto an empty qubit (2), found the cheapest single layer to be
exactly $n+k$ on cycles $m = 4, 5, 6$ with every set of one or two hubs (up to rotation and reflection), adjacent hubs included. Over
$L$ layers with interleaving allowed (a ring edge may be applied for the $(c+1)$-th time once both neighbouring ring edges
have been applied $c$ times), the least total was exactly $L(n+k)$ in every case tried: one hub on $m = 4$ ($L = 2$--$4$), $m = 5$ ($L = 2, 3$) and $m = 6$ ($L = 2$); two non-adjacent hubs on $m = 4$ ($L = 2, 3$) and $m = 5$ ($L = 2$). We have not proved the
bound for interleaved schedules in general.

\begin{lemma}[holes, within their class]
On a cycle $C_m$ with $h$ empty qubits and $n = m-h \ge 3$ nodes, among schedules in which nodes move only onto empty
qubits (never exchange, so their cyclic order is kept), every layer costs at least $n + 2h$.
\end{lemma}
\begin{proof}
Let $q$ be a qubit empty at the start of a layer. If no node enters $q$ during the layer, no node passes it, so with
the cyclic order kept the two nodes beside it stay the same ring-consecutive pair $x$, $x+1$; they are separated by $q$
on one arc and by the other $n-2 \ge 1$ nodes on the other, so they are never adjacent and their ring edge cannot be
applied. Hence each of the $h$ qubits empty at the start is entered during the layer, by distinct operations (each
enters one qubit) that cost 2 and apply no ring edge (a CZ with an empty qubit does nothing); the $n$ ring edges need
$n$ more.
\end{proof}
\noindent The hole schedule meets the bound when no two holes meet (counted at $n = 9, 10, 18$: $15$, $14$, $22$ per
layer); at $n = 8$, where moved holes meet, it uses 4 more over 16 layers. Outside the class the bound fails when holes
are many. The same search, for a single layer from the best start on a bare cycle, finds 6 for $m = 6, 7$ with
$n = 4$ ($n + 2h = 8$, 10), 9 for $m = 8$, $n = 5$ (11) and 11 for $m = 9$, $n = 6$ (12), each by leaving the order (as
the zipper does); with two holes on $m = 8$ it finds $n + 2h = 10$. Where the count rule picks holes on Heron, $h \le 4$; whether $n + 2h$ is least there without the
class restriction is open.

\paragraph{The zipper.} Its count is a search result, not a proof: from every start layout for $n = 4$--$10$ (91,296
start layouts up to symmetry at $n = 10$, of which one reaches the least) and from every minimum-length start at
$n = 12, 14$, no cheaper schedule was found; for odd $n$ the same holds from every minimum-length start for
$n = 5$--$11$ and from every start for $n = 5, 7, 9$ \cite{toqb}.

\section{Signal on noise models}

\paragraph{Measure.} One random classifier of TOQB's ring family (single-qubit rotations per layer, a CZ ring, 16
layers, $\langle Z\rangle$ of qubit 0), 24 random inputs, 256 noise trajectories; the kept signal $\lambda$ is the slope
of the noisy $\langle Z\rangle$ on the noiseless one, $z \approx \lambda z_0 + \delta$. Shots for a fixed precision
scale as $1/\lambda^2$. Two twins: one built from \texttt{ibm\_marrakesh}'s published calibration of 13 October 2026
(the \emph{snapshot}), one from \texttt{ibm\_kingston}'s calibration of 24 October 2026 (the \emph{day twin}), which on
that day had predicted the zipper on hardware (the device's slope on the twin's prediction $1.00$, $0.95$--$1.02$ per
classifier) \cite{paper3}.

\begin{table}[h]
\centering
\small
\begin{tabular}{lcccccc}
\toprule
$n$ & 13 & 14 & 15 & 16 & 17 & 18 \\
\midrule
rule's choice & pendants & pendants & pendants & pendants & pendants & holes \\
two-qubit gates, rule / zipper & 224 / 464 & 256 / 496 & 288 / 544 & 320 / 576 & 352 / 624 & 352 / 656 \\
\midrule
snapshot: rule's $\lambda$ & 0.68--0.73 & 0.64--0.76 & 0.70--0.72 & 0.63--0.73 & 0.61--0.68 & 0.51--0.59 \\
snapshot: zipper's $\lambda$ & 0.09--0.27 & 0.13--0.46 & 0.24--0.31 & 0.17--0.40 & 0.21--0.31 & 0.26--0.34 \\
snapshot: verdict & WIN 4/4 & WIN 4/4 & WIN 4/4 & WIN 4/4 & WIN 4/4 & WIN 4/4 \\
\midrule
day twin: rule's $\lambda$ & 0.76--0.81 & 0.73--0.82 & 0.78--0.81 & 0.77--0.80 & 0.69--0.75 & 0.61--0.68 \\
day twin: zipper's $\lambda$ & 0.64--0.70 & 0.63--0.77 & 0.67--0.69 & 0.60--0.70 & 0.58--0.65 & 0.63--0.66 \\
day twin: ratio & 1.15--1.19 & 1.07--1.18 & 1.15--1.19 & 1.12--1.28 & 1.15--1.28 & 0.98--1.06 \\
day twin: verdict & WIN 4/4 & WIN 4/4 & WIN 4/4 & WIN 4/4 & WIN 4/4 & WIN 3/4 \\
\bottomrule
\end{tabular}
\caption{Declared matches, $L = 16$, four random classifiers each (ranges over the four; on the snapshot both the twin
and a perturbed stand-in were required). WIN: the rule's $\lambda$ larger on at least 3 of 4 classifiers. Both
predictions held.}
\label{tab:match}
\end{table}

\paragraph{Reading.} The direction held on both twins; the size did not. On the snapshot the zipper kept $0.09$--$0.46$
and the rule won by $1.5$--$8.2\times$; on the day twin the zipper kept $0.58$--$0.77$ and the lead was
$1.07$--$1.28\times$ (shots $1.15$--$1.6\times$ fewer), about even at $n=18$. The vanity number (half the two-qubit gates, a shorter duration) holds on
both; the essential number is the day twin's.

\paragraph{Holes against the zipper where holes are many.} On the snapshot, at $L=16$: $n = 8$ (4 holes) the zipper kept
more on every classifier ($0.19$--$0.25$ against $0.10$--$0.19$), $n = 9, 10$ (3, 2 holes) the holes won on every
classifier ($0.26$--$0.37$ and $0.51$--$0.56$ against $0.07$--$0.17$). At $n=8$ the counts tie; the rule's tie goes to
holes and was left so. These are snapshot numbers, so the size of the $n = 9, 10$ leads is an upper bound.

\paragraph{Where the read-out node sits.} Moving the read-out node off the pendants by formula gained $0.01$--$0.05$ at
$n = 14$--$17$ and lost $0.29$--$0.39$ at $n=13$, where the moved node landed next to the cycle's poorest qubit. A
formula that reads no error values cannot avoid such places; it was not adopted.

\paragraph{Error products do not predict the signal.} Two estimators without simulation, the product of $(1-\text{error})$ over
the whole circuit and over the read-out's backward light cone, ranked 64 simulated cells with Spearman correlation
$0.45$--$0.68$ and ordered $82$--$93\,\%$ of clearly different pairs; both failed the test we had fixed for using them
($0.8$ and $90\,\%$ at both depths). The noise reaching the read-out depends on the circuit's shape, not only on how many
errors it contains.

\section{On devices}

\paragraph{Design.} With 18\,s of free device time left, a two-arm comparison was rehearsed on the day twin first: the
pendants' lead of about $0.1$ would have been decided in at most about half the cases even with 1 minute of device
time, because random classifiers at 16 layers spread their noiseless outputs over only $0.08$ (s.d.). We therefore
asked a narrower question: does the pendant schedule keep the signal its day's twin predicts? One arm ($n=13$, $L=16$,
one random classifier, 64 random inputs, 250 shots), compiled and predicted for the Target fetched minutes before,
AGREE if the device slope is within 2 standard errors of the twin's.

\begin{table}[h]
\centering
\small
\begin{tabular}{lccccc}
\toprule
device & twin's $\lambda$ (fixed before) & device $\lambda$ & verdict & billed & zipper's twin $\lambda$ (not run) \\
\midrule
\texttt{ibm\_kingston} & 0.765 & $0.904 \pm 0.124$ & AGREE & 7\,s & 0.613 \\
\texttt{ibm\_fez} & 0.591 & $0.504 \pm 0.102$ & AGREE & 7\,s & 0.432 \\
\bottomrule
\end{tabular}
\caption{The pendant schedule on two devices, 26 October 2026 (224 two-qubit gates, 8--9\,$\mu$s). The device
correlations with the noiseless values ($0.68$, $0.53$) are what 250 shots alone give at this spread ($0.70$, $0.60$).
On \texttt{ibm\_fez} the twin was built from a calibration about 40 minutes older than the one under which the job ran.}
\label{tab:device}
\end{table}

\paragraph{Reading.} On both devices the pendant schedule kept the signal its day's twin predicted, the lower value on
\texttt{ibm\_fez} following that day's lower prediction. These tests rule out a large loss (as Qiskit's circuits showed
on \texttt{ibm\_kingston} on 24 October \cite{paper3}); they cannot see a difference of $0.1$, and no zipper ran on a
device beside them.

\section{Limitations}

The lower bound for pendants is proved for layered schedules in a routing model that excludes measurement,
feed-forward and ancilla-mediated gates, and checked by search only for small interleaved cases; the hole bound holds within a class; the zipper's is a search result. The signal results are on noise models of
two calibrations and one device family, with random rather than trained classifiers, one placement per $n$, and four
classifiers per cell. The device tests are one job each, at a precision of about $0.1$. Count-optimal is not
signal-optimal: duration and where the read-out sits also matter.

\section{Reproducibility}
\sloppy

\begin{itemize}
\item Compiler: PSF-Zero AI front end a26, \path{patches/psf_ai_compile_a26_2026-10-26/} (checks \path{check_a26.py},
\path{test_a26.py}).
\item Runs, each with its declaration fixed by SHA-256 and time before the run (\path{prereg/drafts-2026-10-25/},
\path{prereg/drafts-2026-10-26/}): TOQB \cite{toqb} \path{results/}\ \path{flow-d-2026-10-25}, and, all dated
\path{2026-10-26}, \path{flow-e}, \path{flow-f}, \path{flow-g}, \path{flow-h}, \path{rule-big}, \path{esp-check},
\path{lam-rehearse}, \path{device-d1}, \path{device-d2}, and the searches and cycle counts of Sections 2 and 4, \path{lemma-check} (scripts and outputs; no declaration); the separate review of the pendant code,
\path{pend-review-2026-10-26}.
\item Device jobs: \path{dbfivq0m94ac738tldc0} (\texttt{ibm\_kingston}) and \path{dbfjtvvsstms73fe0jc0}
(\texttt{ibm\_fez}); circuits, manifests and raw counts are in the two device folders.
\end{itemize}

\begin{thebibliography}{99}

\bibitem{psfzero} TN Holdings, ``PSF-Zero: An Analytic Two-Qubit Gate Synthesizer Combined with a Verified,
Ordering-Aware Layout Search for Quantum Circuit Transpilation'', Version 2 (2026).
\url{https://github.com/TN-Holdings-LLC/psf-zero}

\bibitem{paper3} TN Holdings, ``When a Quantum Compiler Looks Like a Better Learner: Confounds in Quantum Machine
Learning Comparisons'', Version 1.2 (2026). \url{https://doi.org/10.5281/zenodo.23630904}

\bibitem{toqb} TN Holdings, TOQB, the test bench for PSF-Zero's training and device runs (2026).
\url{https://github.com/TN-Holdings-LLC/TOQB}

\bibitem{qiskit} A. Javadi-Abhari et al., ``Quantum computing with Qiskit'', arXiv:2405.08810 (2024).
\url{https://arxiv.org/abs/2405.08810}

\bibitem{chamberland2020} C. Chamberland, G. Zhu, T. J. Yoder, J. B. Hertzberg, A. W. Cross, ``Topological and
subsystem codes on low-degree graphs with flag qubits'', Physical Review X 10, 011022 (2020).

\bibitem{perez2020} A. P\'erez-Salinas, A. Cervera-Lierta, E. Gil-Fuster, J. I. Latorre, ``Data re-uploading for a
universal quantum classifier'', Quantum 4, 226 (2020).

\bibitem{vatan2004} F. Vatan, C. Williams, ``Optimal quantum circuits for general two-qubit gates'', Physical Review A
69, 032315 (2004).

\bibitem{shende2004} V. V. Shende, I. L. Markov, S. S. Bullock, ``Minimal universal two-qubit controlled-NOT-based
circuits'', Physical Review A 69, 062321 (2004).

\bibitem{ogorman2019} B. O'Gorman, W. J. Huggins, E. G. Rieffel, K. B. Whaley, ``Generalized swap networks for
near-term quantum computing'', arXiv:1905.05118 (2019).

\end{thebibliography}

\end{document}
